% Working note, not thesis prose -- current state, to revisit together once confirmed.
% Companion to coronary_tree_graph.tex: that file covers why a tree has to be built and what it
% checks out against; this one is the algorithm inside topology/graph.py, build_tree, step by step.

\section*{Building the tree is not a geometric step}

The delivered centerlines need no geometric correction before use. Their points are already
unique and already in physical millimeters (LPS), and no affine transform or spacing correction
touches them, unlike the segmentation volumes, which go through resampling before anything reads
them. There is nothing to reproject and nothing to reslice.

What \texttt{graph.build\_tree(centerline)} does instead is graph construction: turning a file of
separate polylines and flags into one rooted tree, with the connectivity, orientation and labels
that feature extraction needs made explicit per segment.

\begin{itemize}
  \item \textbf{Degree.} \texttt{build\_tree} counts, over every polyline in the file, how many
  times each point index appears as a \texttt{line[0]} or \texttt{line[-1]}. That count is the
  point's degree in the underlying graph.

  \item \textbf{Junctions.} A junction is any point whose degree is not 2, plus every
  \texttt{start\_points} index, pinned in explicitly so an ostium can never be merged straight
  through even if some numerical quirk gave it degree 2.

  \item \textbf{Merging.} \texttt{\_merge\_chains} glues polylines together through the degree-2
  points between junctions, so a segment always runs junction to junction, not polyline to
  polyline as the file stores it. A polyline that forms a closed loop touching no junction at all
  is reported rather than silently dropped; none occur in this dataset, but the check runs
  regardless.

  \item \textbf{Components.} \texttt{\_components}, a union-find over the merged chains, finds
  connected pieces, since 13 of the 1600 sides are not one connected tree.

  \item \textbf{Rooting.} Each component is rooted at its own \texttt{start\_points} index where
  one exists. Two sides (84, 272) have a fragment with no start point at all; that fragment is
  rooted at whichever of its degree-1 nodes lies nearest the already-rooted tree, and the choice
  is logged as a warning rather than made silently.

  \item \textbf{Orientation.} A breadth-first search from each root orients every chain proximal
  to distal. If the search reaches a node it has already visited, that edge closes a cycle, which
  a coronary tree should not have. It is diverted into \texttt{chords} instead of the tree, again
  with a warning. This happens on three sides: 8 left, 455 right, 776 left.

  \item \textbf{Labeling.} A segment's artery label is a majority vote over its interior points,
  not its endpoints. The endpoints are shared junction points and carry the parent vessel's
  label, which would otherwise be read straight into the child segment.

  \item \textbf{Node kind.} Whether a node is an ostium, bifurcation, terminus or pass-through
  comes from how many children it ends up with in the tree that was actually built, not from its
  raw degree in the file. A node whose extra edge was dropped as a chord keeps a raw degree of 3
  but has only one child, so it is typed pass-through, not bifurcation.

  \item \textbf{Vessels.} \texttt{\_build\_vessels} chains same-named segments back together into
  vessels, since a side branch usually cuts a named artery into several segments. Where a name
  splits across more than one child of the same name, the longest downstream run of that name
  continues the vessel; the rest start vessels of their own.
\end{itemize}

None of this changes a single coordinate. It only makes explicit the connectivity, orientation
and labeling the file already encodes but does not spell out on its own.
